% ---- score2 standard preamble (from: python3 score2_tools.py preamble) ----
\documentclass[10pt]{extarticle}
\usepackage[a4paper,margin=1.5cm]{geometry}
\usepackage[T1]{fontenc}
\usepackage{amsmath,amssymb}
\usepackage{graphicx}
\usepackage[export]{adjustbox}
\usepackage{multicol}
\usepackage{enumitem}
\usepackage{booktabs,array,tabularx}
\setlength{\columnsep}{0.7cm}
\setlength{\parindent}{0pt}
\setlength{\parskip}{3pt}
\pagestyle{empty}
\makeatletter
% figure registry (figdefs.tex is generated by the extract step)
\newcommand{\qfigdef}[4]{\@namedef{qf@file@#1}{#2}\@namedef{qf@w@#1}{#3}\@namedef{qf@wide@#1}{#4}}
\newcommand{\qf@check}[1]{\@ifundefined{qf@file@#1}{\PackageError{score2}{Unknown figure id #1}{}}{}}
\newcommand{\qf@img}[4][]{\includegraphics[width=#2\dimexpr\@nameuse{qf@w@#4}pt\relax,max width=#3,max height=0.45\textheight,#1]{\@nameuse{qf@file@#4}}}
% Full-width questions: finalize wraps every question that contains a wide figure in \qfullbegin...\qfullend,
% so the WHOLE question (stem, figures, options) is set across the page in reading order, then the two columns
% resume. The question's head - its number, stem and first figure (\qfullsplit marks the end of the head) -
% is kept on one page, so a number is never left at a page foot with its figure overleaf; the rest may break
% across pages normally, so no big gaps are left. Without a \qfullsplit, the whole question is kept together
% if it is under 45% of a page. Never write these yourself.
\newif\ifqf@full
\newif\ifqf@boxed
\newsavebox{\qf@box}
\newcommand{\qf@open}{\global\qf@fulltrue\par\global\qf@boxedtrue
  \setbox\qf@box\vbox\bgroup\hsize\textwidth\linewidth\textwidth}
\newcommand{\qf@emit}[1]{\par\egroup\global\qf@boxedfalse
  \ifdim\dimexpr\ht\qf@box+\dp\qf@box\relax<#1\textheight\noindent\usebox{\qf@box}\par\else\unvbox\qf@box\fi}
\newcommand{\qfullsplit}{\ifqf@boxed\qf@emit{0.85}\fi}
\newcommand{\qf@close}{\ifqf@boxed\qf@emit{0.45}\fi\global\qf@fullfalse}
\newcommand{\qfullbegin}{\end{multicols}\qf@open}
\newcommand{\qfullend}{\qf@close\begin{multicols}{2}}
\newcommand{\qfullstart}{\qf@open}
\newcommand{\qfulldone}{\qf@close}
% \qfig[scale]{id}: block figure, centred, directly where it is written (size set by the tool).
% Fallback if a wide figure is not inside a full-width question: break out of the columns in place - no floats.
\newcommand{\qfig}[2][1]{\qf@check{#2}%
  \ifqf@full\qf@blk{#1}{#2}\else
  \ifnum\@nameuse{qf@wide@#2}=1 \expandafter\@firstoftwo\else\expandafter\@secondoftwo\fi
  {\end{multicols}\begin{center}\qf@img{#1}{\textwidth}{#2}\end{center}\begin{multicols}{2}}%
  {\qf@blk{#1}{#2}}\fi}
\newcommand{\qf@blk}[2]{\par\smallskip{\centering\qf@img{#1}{\linewidth}{#2}\par}\smallskip}
% \qfigi[scale]{id}: inline figure (inside option lists / tables / minipages); never breaks columns.
\newcommand{\qfigi}[2][1]{\qf@check{#2}\qf@img[valign=c]{#1}{\linewidth}{#2}}
\makeatother
% structure macros - lint reads these to check completeness and figure placement
\newcommand{\papersec}[1]{\par\bigskip{\centering\Large\bfseries PAPER-#1\par}\medskip}
\newcommand{\partsec}[2]{\par\medskip\noindent{\large\bfseries PART-#1: #2}\par\smallskip}
\newcommand{\sect}[1]{\par\smallskip\noindent\textbf{#1}\par}
\newcommand{\qn}[1]{\par\smallskip\noindent\textbf{#1.}~}
\newcommand{\paraq}[2]{\par\smallskip\noindent\textbf{Paragraph for Questions #1 and #2}\par}
\newenvironment{qoptions}{\begin{itemize}[leftmargin=1.8em,itemsep=1pt,topsep=2pt,label={}]}{\end{itemize}}
\newcommand{\opt}[2]{\item[(#1)] #2}
% figure registry (generated by score2_tools.py finalize)
\qfigdef{p1_phy_s1_q03_1}{Score-II - Test 38 (12 Feb 2025)-img/p1_phy_s1_q03_1.png}{119.1}{0}
\qfigdef{p1_phy_s1_q05_1}{Score-II - Test 38 (12 Feb 2025)-img/p1_phy_s1_q05_1.png}{127.8}{0}
\qfigdef{p1_phy_s2_q07_1}{Score-II - Test 38 (12 Feb 2025)-img/p1_phy_s2_q07_1.png}{96.5}{0}
\qfigdef{p1_phy_s2_q12_1}{Score-II - Test 38 (12 Feb 2025)-img/p1_phy_s2_q12_1.png}{57.5}{0}
\qfigdef{p1_phy_s3_q01_1}{Score-II - Test 38 (12 Feb 2025)-img/p1_phy_s3_q01_1.png}{120.4}{0}
\qfigdef{p1_phy_s3_q01_2}{Score-II - Test 38 (12 Feb 2025)-img/p1_phy_s3_q01_2.png}{120.4}{0}
\qfigdef{p1_phy_s3_q03_1}{Score-II - Test 38 (12 Feb 2025)-img/p1_phy_s3_q03_1.png}{179.6}{0}
\qfigdef{p1_phy_s3_q05_1}{Score-II - Test 38 (12 Feb 2025)-img/p1_phy_s3_q05_1.png}{93.2}{0}
\qfigdef{p1_chem_s1_q01_A}{Score-II - Test 38 (12 Feb 2025)-img/p1_chem_s1_q01_A.png}{34.4}{0}
\qfigdef{p1_chem_s1_q01_1}{Score-II - Test 38 (12 Feb 2025)-img/p1_chem_s1_q01_1.png}{41.6}{0}
\qfigdef{p1_chem_s1_q01_7}{Score-II - Test 38 (12 Feb 2025)-img/p1_chem_s1_q01_7.png}{36.0}{0}
\qfigdef{p1_chem_s1_q01_B}{Score-II - Test 38 (12 Feb 2025)-img/p1_chem_s1_q01_B.png}{60.4}{0}
\qfigdef{p1_chem_s1_q01_4}{Score-II - Test 38 (12 Feb 2025)-img/p1_chem_s1_q01_4.png}{44.5}{0}
\qfigdef{p1_chem_s1_q01_8}{Score-II - Test 38 (12 Feb 2025)-img/p1_chem_s1_q01_8.png}{36.2}{0}
\qfigdef{p1_chem_s1_q01_5}{Score-II - Test 38 (12 Feb 2025)-img/p1_chem_s1_q01_5.png}{115.4}{0}
\qfigdef{p1_chem_s1_q01_6}{Score-II - Test 38 (12 Feb 2025)-img/p1_chem_s1_q01_6.png}{69.6}{0}
\qfigdef{p1_chem_s1_q02_1}{Score-II - Test 38 (12 Feb 2025)-img/p1_chem_s1_q02_1.png}{105.0}{0}
\qfigdef{p1_chem_s1_q02_2}{Score-II - Test 38 (12 Feb 2025)-img/p1_chem_s1_q02_2.png}{45.5}{0}
\qfigdef{p1_chem_s1_q02_B}{Score-II - Test 38 (12 Feb 2025)-img/p1_chem_s1_q02_B.png}{120.7}{0}
\qfigdef{p1_chem_s1_q02_4}{Score-II - Test 38 (12 Feb 2025)-img/p1_chem_s1_q02_4.png}{150.9}{0}
\qfigdef{p1_chem_s1_q02_5}{Score-II - Test 38 (12 Feb 2025)-img/p1_chem_s1_q02_5.png}{93.8}{0}
\qfigdef{p1_chem_s2_q07_1}{Score-II - Test 38 (12 Feb 2025)-img/p1_chem_s2_q07_1.png}{171.0}{0}
\qfigdef{p1_chem_s2_q07_2}{Score-II - Test 38 (12 Feb 2025)-img/p1_chem_s2_q07_2.png}{207.5}{0}
\qfigdef{p1_chem_s2_q07_3}{Score-II - Test 38 (12 Feb 2025)-img/p1_chem_s2_q07_3.png}{125.3}{0}
\qfigdef{p1_chem_s2_q07_4}{Score-II - Test 38 (12 Feb 2025)-img/p1_chem_s2_q07_4.png}{44.1}{0}
\qfigdef{p1_chem_s2_q08_A}{Score-II - Test 38 (12 Feb 2025)-img/p1_chem_s2_q08_A.png}{167.5}{0}
\qfigdef{p1_chem_s2_q08_B}{Score-II - Test 38 (12 Feb 2025)-img/p1_chem_s2_q08_B.png}{110.4}{0}
\qfigdef{p1_chem_s2_q08_C}{Score-II - Test 38 (12 Feb 2025)-img/p1_chem_s2_q08_C.png}{108.2}{0}
\qfigdef{p1_chem_s2_q08_D}{Score-II - Test 38 (12 Feb 2025)-img/p1_chem_s2_q08_D.png}{104.5}{0}
\qfigdef{p1_chem_s2_q10_1}{Score-II - Test 38 (12 Feb 2025)-img/p1_chem_s2_q10_1.png}{203.5}{0}
\qfigdef{p1_chem_s3_q01_1}{Score-II - Test 38 (12 Feb 2025)-img/p1_chem_s3_q01_1.png}{245.2}{0}
\qfigdef{p1_chem_s3_q02_1}{Score-II - Test 38 (12 Feb 2025)-img/p1_chem_s3_q02_1.png}{74.9}{0}
\raggedcolumns
\makeatletter\@beginparpenalty=10000\makeatother
\begin{document}
\begin{center}{\Large\bfseries JEE (Main+Advanced) Score-II --- Test 38 (12 Feb 2025)}\end{center}
\begin{multicols}{2}

\partsec{1}{PHYSICS}
\sect{Section-I(i)}

\qn{1} Four students Ram, Shyam, Ghanshyam and Radheshyam perform an optical bench experiment to find focal length of a convex lens. Ram forms image of a distant luminous object on a small screen mounted on the optical bench. Other 3 students perform usual optical bench experiment. Their position of object needle, lens and image needle are given below. Which of the student is most accurate in his result?
\begin{center}\small
\begin{tabular}{|c|c|c|c|}
\hline
\textbf{Name} & \textbf{Object} & \textbf{Lens} & \textbf{Image} \\
\hline
Ram & $-$ & 15.0 cm & 30.0 cm \\
\hline
Shyam & 15.0 cm & 75.0 cm & 95.0 cm \\
\hline
Ghanshyam & 15.0 cm & 45.0 cm & 75.0 cm \\
\hline
Radheshyam & 15.0 cm & 60.0 cm & 82.5 cm \\
\hline
\end{tabular}
\end{center}
\begin{qoptions}
\opt{A}{Ram}
\opt{B}{Ghanshyam}
\opt{C}{Radheshyam}
\opt{D}{All are equally accurate}
\end{qoptions}

\qn{2} If the third harmonic vibration in an open air pipe equals the fifth harmonic of vibration in a closed air pipe, the possible values of the length of air column in the closed pipe and open pipe, respectively, are (neglect end corrections) :-
\begin{qoptions}
\opt{A}{240 cm, 120 cm}
\opt{B}{60 cm, 72 cm}
\opt{C}{200 cm, 120 cm}
\opt{D}{180 cm, 150 cm}
\end{qoptions}

\qn{3} A ring of mass M and radius R lies in x-y plane with its centre at origin as shown. The mass distribution of ring is non-uniform such that at any point P on the ring, the mass per unit length is given by $k=k_0\left(\cos^2\theta+\dfrac\theta\pi\right)$ (where $k_0$ is a positive constant, $\theta$ is in radians). Then the moment of inertia of the ring about z-axis is
\qfig{p1_phy_s1_q03_1}
\begin{qoptions}
\opt{A}{$MR^2$}
\opt{B}{$\dfrac12MR^2$}
\opt{C}{$\dfrac12\dfrac{M^2}{k_0}R$}
\opt{D}{$\dfrac1\pi\dfrac{M^2}{k_0}R$}
\end{qoptions}

\qn{4} An $\alpha$ particle is moving along a circle of radius R with a constant angular velocity. Point A lies in the same plane at a distance 2R from the centre. The point A records magnetic field produced by the $\alpha$ particle. If the minimum time interval between two successive times at which A records zero magnetic field is `t', find frequency of revolution of the $\alpha$-particle.
\begin{qoptions}
\opt{A}{$\dfrac1t$}
\opt{B}{$\dfrac1{3t}$}
\opt{C}{$\dfrac1{6t}$}
\opt{D}{$\dfrac1{2t}$}
\end{qoptions}

\qn{5} Consider two coherent monochromatic (wavelength $\lambda$) sources $S_1$ and $S_2$ separated by distance d. The ratio of intensity of $S_1$ and that of $S_2$ at point P, is 4. The distance of point P from $S_1$ is (if the resultant intensity at point P is equal to $\dfrac94$ times of intensity of $S_1$)
\qfig{p1_phy_s1_q05_1}
\relax[Given $\angle S_2S_1P$ is $90^\circ$] (d $>$ 0, n is a positive integer)
\begin{qoptions}
\opt{A}{$\dfrac{d^2-n^2\lambda^2}{2n\lambda}$}
\opt{B}{$\dfrac{d^2+n^2\lambda^2}{2n\lambda}$}
\opt{C}{$\dfrac{n\lambda d}{\sqrt{d^2-n^2\lambda^2}}$}
\opt{D}{$\dfrac{2n\lambda d}{\sqrt{d^2-n^2\lambda^2}}$}
\end{qoptions}

\qn{6} A conducting disc of conductivity $\sigma$ has radius a and thickness $\delta$. A uniform magnetic field $\vec B$ is applied in a direction perpendicular to the plane of the disc. If the magnetic field changes with time at the rate of R, then find the power dissipated in the disc due to the induced current.
\begin{qoptions}
\opt{A}{$\dfrac{\pi\delta\sigma a^4}8(R)^2$}
\opt{B}{$\dfrac{\pi\delta\sigma a^4}{12}(R)^2$}
\opt{C}{$\dfrac{\pi\delta\sigma a^4}4(R)^2$}
\opt{D}{$\dfrac{\pi\delta\sigma a^4}6(R)^2$}
\end{qoptions}

\sect{Section-I(ii)}

\qn{7} The system in figure lies on a frictionless horizontal table. Both masses are m, and the two spring constant are nK (where n is a numerical factor) and K as shown in figure. The springs have the same relaxed length. Assuming that it is possible to set up initial conditions so that the masses oscillate back and forth with the two springs always having equal lengths at any given instant. Choose CORRECT statement(s) -
\qfig{p1_phy_s2_q07_1}
\begin{qoptions}
\opt{A}{Relative to mean position, if displacement of left mass is x, then displacement of right mass is 2x.}
\opt{B}{Relative to mean position, if displacement of left mass is x, then displacement of right mass is x.}
\opt{C}{Value of n is $\dfrac32$.}
\opt{D}{Value of n is $\dfrac23$.}
\end{qoptions}

\qn{8} In a new system of units, the unit of mass is 1000 kg (1 metric ton), unit of length is 1000m (1km) and unit of time is 3600 s(1hr). Select correct statement (s).
\begin{qoptions}
\opt{A}{The numerical value of 500 m in this new system is 0.5.}
\opt{B}{The numerical value of 7200 s in this new system is 2.}
\opt{C}{The numerical value of $\dfrac1{36}$J in this new system is $3.6\times10^{-4}$.}
\opt{D}{The numerical value of $\dfrac1{36}$N in this new system is 0.36.}
\end{qoptions}

\qn{9} Consider a hypothetical hydrogen like atom. The wavelength of spectral line from $p^{\text{th}}$ orbit to ground state is given by $\lambda=\dfrac{3500p^2}{p^2-1}$ in \AA. Choose the correct option(s) for considered hypothetical atom. (hC =1240 eV-nm)
\begin{qoptions}
\opt{A}{The ionization potential of considered hypothetical atom is 3.54 V approximately}
\opt{B}{The ionization potential of considered hypothetical atom is 0.354 V approximately}
\opt{C}{The wavelength of least energetic photon in the spectral line series is 466.67 nm approximately}
\opt{D}{The wavelength of maximum energetic photon in the spectral line series is 350.00 nm approximately}
\end{qoptions}

\qn{10} The potential energy function of a particle is given by $U(r)=\dfrac A{2r^2}-\dfrac B{3r}$, where A and B are constants and r is the radial distance from the centre of the force. Choose the correct option(s).
\begin{qoptions}
\opt{A}{The equilibrium distance will be $r_0=\dfrac{2A}B$}
\opt{B}{The equilibrium distance will be $r_0=\dfrac{3A}B$}
\opt{C}{If the total energy of the particle is $\dfrac{B^2}{6A}$, then its radial velocity will vanish at $\dfrac{r_0}3$, where $r_0$ is the equilibrium distance}
\opt{D}{If the total energy of the particle is $\dfrac{B^2}{6A}$, then its radial velocity will vanish at $\dfrac{r_0}2$, where $r_0$ is the equilibrium distance}
\end{qoptions}

\qn{11} Radon decays to polonium according to the decay scheme given by ${}^{222}_{88}\text{Rn}\xrightarrow[3.8\text{ days}]{\alpha-\text{decay}}{}^{218}_{84}\text{Po}$. State which of the following statement(s) will be true in respect of this decay : (Given : $\ln2=0.693$)
\begin{qoptions}
\opt{A}{If the original mass of radon in a closed container was 1mg, the amount of radon left after 11.4 days will be 0.125 mg.}
\opt{B}{Activity of radon after 7.6 days will be $\dfrac{A_0}4$ where $A_0$ is original activity of sample at t = 0.}
\opt{C}{If the original number of radon in a closed container was 8 atoms, the number of radon left after 7.6 day will be 2 atoms.}
\opt{D}{$P_0$ is more stable than Rn.}
\end{qoptions}

\qn{12} Three processes compose a thermodynamic cycle shown in the accompanying PV diagram of an ideal gas.
\qfig{p1_phy_s2_q12_1}
Process $1\to2$ takes place at constant temperature (300 K). During this process 60 J of heat enters the system.\\
Process $2\to3$ takes place at constant volume. During this process 40 J of heat leaves the system.\\
Process $3\to1$ is adiabatic. $T_3$ is 275 K.
\begin{qoptions}
\opt{A}{The change in internal energy of the system during process $3\to1$ is +40 J}
\opt{B}{The change in internal energy of the system during process $3\to1$ is 0}
\opt{C}{The change in internal energy of the system during cyclic process is 40 J}
\opt{D}{The change in internal energy of the system during cyclic process is 0}
\end{qoptions}

\sect{Section-II}

\qn{1} \qfigi[1]{p1_phy_s3_q01_1}\\
In the given circuit the resistance R can be varied. Battery having emf 18 volt is having internal resistance 3 $\Omega$. The value of R is selected such that power dissipated in it is maximum. Now 6 $\Omega$ resistance is replaced by an inductor having inductance L = 6mH so circuit will appear like the circuit shown in diagram given below\\
\qfigi{p1_phy_s3_q01_2}\\
Find time constant for above circuit in milliseconds.

\qn{2} An ac ammeter is used to measure current in the circuit. When a given direct current passes through the circuit, the ac ammeter reads 3 amp. When another alternating current passes through the circuit, the ac ammeter reads 4 amp. Then the reading of this ammeter (in Ampere) if both dc and ac flows through the circuit simultaneously is \underline{\hspace{1cm}}.

\qn{3} In the circuit shown length of AB is 100 cm. Then find the value of $\ell_1$ (in cm). Where $\ell_1$ corresponds to null point.
\qfig{p1_phy_s3_q03_1}

\qn{4} In a meter bridge experiment, the unknown resistance was connected in left branch and a resistance box in right branch. The resistance in resistance box is 50$\Omega$ (known acurately). There is no end correction. The least count of meter bridge is 1mm. The null point was obtained at 40.0cm from left end. Now the radius of wire was measured by a screw guage with pitch of 1mm and no. of divisions on circular scale as 100. The screw guage read 1 main scale division and 60 circular scale divisions. The length of resistor was measured using a vernier callipers with 1mm as 1 main scale division and 9 main scale division = 10 vernier scale divisions. The length is 2.2 cm as main scale reading and $5^{\text{th}}$ vernier scale division coinciding. The percentage error in measurement of resistivity is.

\qn{5} The figure shows a RC circuit with a parallel plate capacitor. Before switching on the circuit, plate A of the capacitor has a charge $-Q_0$ while plate B has no net charge. Now, at t = 0, the circuit is switched on. How much time (in second) will elapse before the net charge on plate A becomes zero.\\
(Given C = 1$\mu$F, $Q_0$ = 1mC, $\epsilon$ = 1000 V and R $=\dfrac{2\times10^6}{\ln3}\Omega$)
\qfig{p1_phy_s3_q05_1}

\qn{6} A rod of length L and mass M is located along the x-axis with ends at x = 0 and x = L. This rod has a special property such that if we make a cut at any arbitrary x = a coordinate and consider the rod between x = 0 and x = a, the centre of mass of this portion of the rod is located at a/4. If the linear mass density $\lambda(x)$ of the rod as a function of x is given as $\lambda(x)=\dfrac{kM}{nx^{1/2n}L^{1/4n}}$. Find n.

\partsec{2}{CHEMISTRY}
\sect{Section-I(i)}

\qn{1} Which of the following statement is \textbf{INCORRECT} ?
\begin{qoptions}
\opt{A}{\qfigi{p1_chem_s1_q01_A} $\xrightarrow[\text{(ii) H}^+]{\text{(i) P}}$ \qfigi{p1_chem_s1_q01_1}\\ In above reaction compound P is \qfigi{p1_chem_s1_q01_7}}
\opt{B}{\qfigi{p1_chem_s1_q01_B}\\ \qfigi{p1_chem_s1_q01_4}\\ In above reaction compound Q = \qfigi{p1_chem_s1_q01_8}}
\opt{C}{Order of rate of nucleophilic addition reaction is\\ \qfigi{p1_chem_s1_q01_5}}
\opt{D}{IUPAC name of compound \qfigi{p1_chem_s1_q01_6}\\ 3-Amino-2-formyl cyclopropane carboxylic acid}
\end{qoptions}

\qn{2} Which of the following is \textbf{CORRECT} ?
\begin{qoptions}
\opt{A}{In the reaction\\ \qfigi{p1_chem_s1_q02_1}\\ compound \qfigi{p1_chem_s1_q02_2} is obtained as major product}
\opt{B}{\qfigi{p1_chem_s1_q02_B}}
\opt{C}{Correct order of stability is\\ \qfigi{p1_chem_s1_q02_4}}
\opt{D}{Stability of\\ \qfigi{p1_chem_s1_q02_5}}
\end{qoptions}

\qn{3} Which of the following statement/relationship is incorrect?
\begin{qoptions}
\opt{A}{Reducing power in aq. medium : Li $>$ K $>$ Na}
\opt{B}{Solubility in water : Ca(OH)$_2$ $<$ Sr(OH)$_2$ $<$ Ba(OH)$_2$}
\opt{C}{NO is obtained as one of the product when Sn reacts with conc. HNO$_3$}
\opt{D}{Lithium carbide produces acetylene gas on hydrolysis whereas aluminium carbide produces methane gas on hydrolysis}
\end{qoptions}

\qn{4} Considering complete hydrolysis of XeF$_2$ and XeF$_4$,\\
$\text{XeF}_2\xrightarrow{\text{H}_2\text{O}}(\text{P})+(\text{Q})+(\text{R})$\\
$\text{XeF}_4\xrightarrow{\text{H}_2\text{O}}(\text{P})+(\text{Q})+(\text{R})+(\text{S})$\\
Then shape of compound `S' is
\begin{qoptions}
\opt{A}{Linear}
\opt{B}{Trigonal planar}
\opt{C}{Trigonal pyramidal}
\opt{D}{T-shape}
\end{qoptions}

\qn{5} Which of the following sols can not be coagulated by Al$^{3+}$ ion
\begin{qoptions}
\opt{A}{Sols of starch}
\opt{B}{As$_2$S$_3$ sol}
\opt{C}{Silver sol}
\opt{D}{Methylene blue sol}
\end{qoptions}

\qn{6} If a certain amount of ideal gas in \textbf{state-1} (2L, 5 atm, 300 K) is subjected to expand reversibly and isothermally to \textbf{state-2} (V = 10L) and then the gas is restored to its initial \textbf{state-1} by compressing it using an external pressure of $P_{\text{ext.}}$ in single step, then the correct information(s) is/are :\\
(Given : log2 = 0.3, log3 = 0.48, log5 = 0.7)
\begin{qoptions}
\opt{A}{Work done by the gas during expansion from state 1 to State 2 is 8.06 L-atm (approx.)}
\opt{B}{Heat exchanged in the 2nd step (State 2 to State 1) is $-40$ J (approx.)}
\opt{C}{Net heat gained by the surrounding is 23.88 L-atm (approx.)}
\opt{D}{Net change in entropy of the system is positive}
\end{qoptions}

\sect{Section-I(ii)}

\qn{7} Which of the following statement(s) is/are \textbf{NOT INCORRECT} ?
\begin{qoptions}
\opt{A}{compound\\ \qfigi{p1_chem_s2_q07_1}}
\opt{B}{Correct order of basic strength\\ \qfigi{p1_chem_s2_q07_2}}
\opt{C}{Order of acidic strength\\ \qfigi{p1_chem_s2_q07_3}}
\opt{D}{Alkene is obtained during \textbf{Kolbe's electrolysis} of compound\\ \qfigi{p1_chem_s2_q07_4}}
\end{qoptions}

\qn{8} Select reaction(s) with \textbf{INCORRECT} major product
\begin{qoptions}
\opt{A}{\qfigi{p1_chem_s2_q08_A}}
\opt{B}{\qfigi{p1_chem_s2_q08_B}}
\opt{C}{\qfigi{p1_chem_s2_q08_C}}
\opt{D}{\qfigi{p1_chem_s2_q08_D}}
\end{qoptions}

\qn{9} Select correct statement(s)/relationship
\begin{qoptions}
\opt{A}{On strong heating NH$_4$NO$_2$ produces paramagnetic gaseous substance}
\opt{B}{Order of thermal stability : LiNO$_3$ $<$ NaNO$_3$ $<$ KNO$_3$}
\opt{C}{SiO$_2$ is non reactive with NaOH(aq.) whereas it reacts with NaOH($\ell$) on fusion}
\opt{D}{Fe and Zn are commercially extracted by carbon reduction method}
\end{qoptions}

\qn{10} Based on the given Ellingham diagram showing the change in free energy ($\Delta$G$^\circ$ in kJ/mole of O$_2$) for oxide formation vs temperature, select the correct statement(s)
\qfig{p1_chem_s2_q10_1}
\begin{qoptions}
\opt{A}{At 1100$^\circ$C, CO$_2$ + C $\rightarrow$ 2CO is favourable whereas Al$_2$O$_3$ + 3C $\rightarrow$ 2Al + 3CO is unfavourable}
\opt{B}{Al can reduce MgO below 1500$^\circ$C}
\opt{C}{Mg can reduce ZrO$_2$ but Al can not reduce ZrO$_2$ below 1100 $^\circ$C}
\opt{D}{CO is better reducing agent than coke at relatively lower (below 800 $^\circ$C) temperatures for reduction of FeO}
\end{qoptions}

\qn{11} For an electrochemical cell\\
A(s)$|$A$^{x+}$(aq. , 1M)$||$B$^{2x+}$(aq. , 0.4M)$|$B(s),\\
$E_{\text{cell}}=0$ and the value of $\Delta_rG^\circ$ is half of $\Delta_rH^\circ$ at 300 K. Considering log 2 = 0.3, F = 96500 C mol$^{-1}$ and R = 8.3 JK$^{-1}$mol$^{-1}$. Select the correct relationship(s)
\begin{qoptions}
\opt{A}{$E^\circ_{\text{cell}}$ is positive}
\opt{B}{$\dfrac{dE^\circ_{\text{cell}}}{dT}$ is negative}
\opt{C}{$\Delta_rS^\circ\simeq-7.65$ JK$^{-1}$mol$^{-1}$}
\opt{D}{Cell reaction in standard condition and at 300 K is endothermic in nature}
\end{qoptions}

\qn{12} A 0.252 g sample of monoprotic organic base (amine) is dissolved in water and titrated against 0.07 M H$_2$SO$_4$ solution at 25$^\circ$C. After addition of 20 ml of acid, a pH of 10.7 is recorded. Equivalence point is observed when 40 ml of H$_2$SO$_4$ is added. (log 2 = 0.3, log 3 = 0.48, log 5 = 0.7)\\
Based on the above information select correct statement(s)
\begin{qoptions}
\opt{A}{Molar mass of the organic base is 45 g/mole}
\opt{B}{Resultant solution after addition of 20 ml of H$_2$SO$_4$ can act as basic buffer solution.}
\opt{C}{$K_b$ of the organic base is $5\times10^{-4}$}
\opt{D}{Solution formed at equivalence point is basic in nature.}
\end{qoptions}

\sect{Section-II}

\qn{1} (i) Consider following set of reactions
\qfig{p1_chem_s3_q01_1}
Degree of unsaturation of compound B is `\textbf{y}'\\
\hspace*{1.5cm}\&\\
\textbf{(ii)} During detection of sulphur in \textbf{Lassaigne's test,} violet coloured compound with formula ``Na$_x\,[\text{Fe(CN)}_5\text{NOS}]$ '' is formed, then what is the value of $\dfrac yx$.

\qn{2}
\qfig{p1_chem_s3_q02_1}
In the above given compound\\
\textbf{X} = Total number of functional groups that can be reduced by LAH (Lithium aluminium hydride),\\
\textbf{Y} = Total number of functional groups that can be reduced by SBH (sodium brohydride)\\
\hspace*{1.5cm}\&\\
\textbf{Z} = Total number of isomeric alcohols of molecular formula C$_5$H$_{12}$O (including stereoisomers) which can give blue colour in \textbf{Victor - Meyer test}\\
Then value of $\left(\dfrac{X+Y}Z\right)$ = ...........

\qn{3} Find the sum of spin only magnetic moment (in B.M) of the complexes (A, B, C) formed in following processes.\\
(i) MnSO$_4$(aq.) + Pb$_3$O$_4$ $\xrightarrow[\Delta]{\text{H}^+}$ complex ion (A)\\
(ii) NiCl$_2$(aq.) + $\underset{\text{(Excess)}}{\text{ethylenediamine}}$ $\rightarrow$ complex ion (B)\\
(iii) CuSO$_4$(aq.) + $\underset{\text{(Excess)}}{\text{KCN}}$ $\rightarrow$ complex ion (C)

\qn{4} Among the following reactions, how many reactions produce precipitate of basic salt?\\
(i) Mg$^{+2}$(aq.) $\xrightarrow{\text{Na}_2\text{CO}_3\text{ (aq.)}}$\\
(ii) Zn$^{+2}$(aq.) $\xrightarrow{\text{Na}_2\text{CO}_3\text{ (aq.)}}$\\
(iii) Pb$^{+2}$(aq.) $\xrightarrow{\text{Na}_2\text{CO}_3\text{ (aq.)}}$\\
(iv) Hg$^{+2}$(aq.) $\xrightarrow{\text{Na}_2\text{CO}_3\text{ (aq.)}}$\\
(v) Cu$^{2+}$(aq.) $\xrightarrow{\text{Na}_2\text{CO}_3\text{ (aq.)}}$\\
(vi) Zn$^{2+}$(aq.) $\xrightarrow{\text{NaHCO}_3\text{ (aq.)}}$\\
(vii) Ca$^{2+}$(aq.) $\xrightarrow{\text{Na}_2\text{CO}_3\text{ (aq.)}}$\\
(viii) Hg$_2^{2+}$(aq.) $\xrightarrow{\text{Na}_2\text{CO}_3\text{ (aq.)}}$\\
(ix) Hg$_2^{2+}$(aq.) + NO$_3^-$(aq.) $\xrightarrow{\text{NH}_4\text{OH (aq.)}}$\\
(x) Cr$^{3+}$ + (aq.) $\xrightarrow{\text{Na}_2\text{CO}_3\text{ (aq.)}}$

\qn{5} 2.4g of pure carbon is burnt in limited supply of O$_2$ to produce CO and CO$_2$(g). After 100\% conversion of carbon into CO and CO$_2$, the gaseous mixture is treated with I$_2$O$_5$ for the estimation of CO. If liberated I$_2$ in this process require 300ml of 0.1N hypo solution for complete titration then \% by mole of CO in the gaseous mixture of CO and CO$_2$ is \underline{\hspace{1cm}}

\qn{6} 1 m, 1.06 kg of urea solution is kept at temperature 10$^\circ$C below its freezing point. It is observed that 50\% of solvent gets separated. Find the mass (in gm) of solvent freezes out if the same solution is kept at temperature 15$^\circ$C below its freezing point.

\partsec{3}{MATHEMATICS}
\sect{Section-I(i)}

\qn{1} Let $f(x)=2x-2x^2$, $x\in[0,1]$. Let $f_n(x)=\underbrace{fofofo\dots\dots\dots\dots of(x)}_{n\text{ times}}$ such that $\displaystyle\int_0^1f_{2017}(x)dx=\dfrac{2^p}{2^q+1}$, then value of $q-p$ is
\begin{qoptions}
\opt{A}{1}
\opt{B}{2}
\opt{C}{3}
\opt{D}{4}
\end{qoptions}

\qn{2} Let $g:R\to R$ be a differentiable function $(g'(0)\neq0)$, and $f:R\to R$ be a continuous function, such that\\
$g(x+y)=g(x)+g(y)+1\ \forall\ x,y\in R$ and\\
$g\left(f^2(x)+2+x^2f^2(x)\right)-g\left(x^2f(x)+3f(x)\right)=0,\ \forall\ x\in R$.\\
and $f(2023)=f(2024)\neq f(0)$, then the value of $\left(f\left(\dfrac12\right)+f(2)\right)$ can be equal to
\begin{qoptions}
\opt{A}{2}
\opt{B}{3}
\opt{C}{$\dfrac{13}5$}
\opt{D}{$\dfrac85$}
\end{qoptions}

\qn{3} If\\
{\footnotesize $\begin{array}{l}\vec P=(\cos x)\hat i+(\sin x)\hat j+(\cos2x)\hat k,\\ \vec Q=(\tan x)\hat i+(\sin3x)\hat j+(\cos4x)\hat k\text{ and}\\ \vec R=(\cos3x)\hat i+(\sin5x)\hat j+(\cos6x)\hat k\end{array}$}\\
are three variable vectors given, then\\ $\displaystyle\int_0^\pi\left(\left(\vec P\times\vec Q\right).\vec R\right)dx$ belongs to -
\begin{qoptions}
\opt{A}{$(-1,0)$}
\opt{B}{$(0,1)$}
\opt{C}{$(-\infty,-1)\cup(1,\infty)$}
\opt{D}{$(-1,1)$}
\end{qoptions}

\qn{4} Let volume of a tetrahedron ABCD is $\dfrac{81}2$ cube unit and volume of parallelopiped whose three coterminous edges are line segments joining centroid of any face of tetrahedron with centroids of its other three faces is V cube unit, then V is
\begin{qoptions}
\opt{A}{9}
\opt{B}{8}
\opt{C}{7}
\opt{D}{3}
\end{qoptions}

\qn{5} If $f(x)=[2+5|n|\sin x]$ where $n\in I$, $x\in[0,\pi]$ has exactly 19 points of non differentiability then the possible value of $|n|$ is (where $[\,*\,]$ represents G.I.F.)
\begin{qoptions}
\opt{A}{1}
\opt{B}{2}
\opt{C}{3}
\opt{D}{4}
\end{qoptions}

\qn{6} Number of ways in which 4A, 1B, 1C and 1D can be distributed among 4 persons such that no one gets all 4A and every one gets at least one letter is
\begin{qoptions}
\opt{A}{832}
\opt{B}{2432}
\opt{C}{2077}
\opt{D}{608}
\end{qoptions}

\sect{Section-I(ii)}

\qn{7} Let A(0, 6, 8) and B(15, 20, 0) are two given points and $P(\lambda,0,0)$ is a point on x-axis such that PA + PB is minimum. If image of origin along plane mirror passing through P, A, B is $(\alpha,\beta,\gamma)$ then
\begin{qoptions}
\opt{A}{$9(\alpha+\beta)=20$}
\opt{B}{$9(\beta+\gamma)=20$}
\opt{C}{$9(\gamma+\alpha)=20$}
\opt{D}{$9(\alpha+\beta+\gamma)=60$}
\end{qoptions}

\qn{8} Let F and G be two real-valued functions defined on R (the set of all real numbers).\\
Consider $F(x)=\begin{cases}|x+1|&,\ x\leqslant0\\ x&,\ x>0\end{cases}$, $G(x)=\begin{cases}|x|+1&,\ x\leqslant0\\ -|x-2|&,\ x>0\end{cases}$ and $H(x)=F(x)+G(x)$.\\
Points of discontinuity of H(x) are
\begin{qoptions}
\opt{A}{x = 0}
\opt{B}{x = 1}
\opt{C}{x = 2}
\opt{D}{x = $-1$}
\end{qoptions}

\qn{9} A variable line L intersects the parabola $x^2=y$ at points A and B whose x-coordinates are $\alpha$ and $\beta$ with $\alpha<\beta$. The area of figure enclosed by the segment AB and the parabola is always equal to $\dfrac43$. The variable segment AB has its middle point as M. Then the locus of mid point M is a parabola $f(x,y)=0$, whose focus is S. If from point $P(-2,0)$ two tangents are drawn to the parabola $f(x,y)=0$ touching it at the point Q and R, then 16 (SQ) (SR) is always less than
\begin{qoptions}
\opt{A}{89}
\opt{B}{91}
\opt{C}{98}
\opt{D}{101}
\end{qoptions}

\qn{10} If A, B, C are the angles of an acute angled triangle ABC and\\
\resizebox{\linewidth}{!}{$D=\begin{vmatrix}(\tan B+\tan C)^2&\tan^2A&\tan^2A\\ \tan^2B&(\tan A+\tan C)^2&\tan^2B\\ \tan^2C&\tan^2C&(\tan A+\tan B)^2\end{vmatrix}$,}\\
then the value of $\dfrac D{1000}$ can be
\begin{qoptions}
\opt{A}{$-1$}
\opt{B}{0}
\opt{C}{2}
\opt{D}{3}
\end{qoptions}

\qn{11} Let a, b, c are non-zero real numbers such that $\dfrac1a,\dfrac1b,\dfrac1c$ are in arithmetic progression and a, b, $-2$c are in geometric progression, then which of the following statement(s) is(are) must be true ?
\begin{qoptions}
\opt{A}{$a^2,b^2,4c^2$ are in geometric progression}
\opt{B}{$-2a,b,-2c$ are in arithmetic progression}
\opt{C}{$a^3+b^3+c^3-3abc=0$}
\opt{D}{$a^2,b^2,c^2$ are in harmonic progression}
\end{qoptions}

\qn{12} Let $z=a+ib$ (where a, b, $\in$ R) and $i=\sqrt{-1}$ such that $|2z+3i|=|z^2|$. Identify the correct statement(s)?
\begin{qoptions}
\opt{A}{$|Z|_{\text{maximum}}$ is equal to 3.}
\opt{B}{$|Z|_{\text{minimum}}$ is equal to 1}
\opt{C}{If $|Z|_{\text{maximum}}$ when $z=\alpha+i\beta$ ($\alpha,\beta\in$ R and $i=\sqrt{-1}$) then $(\alpha^3+\beta^3)$ is equal to 27}
\opt{D}{If $|Z|_{\text{minimum}}$ when $z=x+iy$ ($\alpha,\beta\in$ R and $i=\sqrt{-1}$) then $(x^2+2y^2)$ is equal to 2}
\end{qoptions}

\sect{Section-II}

\qn{1} Let $f:R\to R$ be a function defined as\\
$f(x+y)=f(x)+f(y)-3xy(x+y)-1\ \forall\ x,y\in R$, $f'(0)=-1$.\\
Find the least positive integer value of x which satisfies the inequality\\
$1-f(x)-f^3(x)>f(1-5x)$.

\qn{2} A number is chosen at random from the numbers 10 to 99. By seeing the number a man will laugh if product of the digits is 12. If he choose three numbers with replacement then the probability that he will laugh at least once is $1-\left(\dfrac pq\right)^3$. Where (p, q are co-prime) then the value of $\left|\dfrac3{q-p}\right|$ is

\qn{3} If $f(x)+f(y)=\dfrac1x+\dfrac1y\ \forall x,y\in R-\{0\}$ and $\displaystyle\int_2^3\dfrac{3(f(x))^5-f(x)}{1-(f(x))^4}dx=\dfrac12\log\dfrac{2^\alpha}{3^\beta}$, then $\dfrac\beta\alpha=$

\qn{4} Let $f(x)=x^3-3x^2+4x-1$, then the value of $\displaystyle\int_{4/27}^{50/27}f(f(x))dx$ is

\qn{5} If $p=\displaystyle\sum_{r=1}^{50}\left\{\sum_{k=1}^{r}\left(\dfrac{(-1)^{r-1}}k.{}^{50}C_r\right)\right\}$ then the value of $100(p)$ equals

\qn{6} If z is a complex number such that $|z|=1$, then the maximum value of $\dfrac{\left|z^4+z^3-2z^2i+z+1\right|^2}5$ is

\end{multicols}
\end{document}
